\subsection{Declarations}
\label{ssec:declarations}

Declarations, of values and functions, may be made either at the top level of
a script, or within a local block.  Standard scoping rules apply. 

A \emph{value declaration} is written using the keyword |val|.
Figure~\ref{fig:example-script} contains several examples.  Value declarations
may pattern match against tuples.  For example
%
\begin{scala}
val (x,y) = pair
\end{scala}
%
This is particularly useful to unpack a tuple returned by a function. 

A \emph{function declaration} is written using the keyword |def|.  Again,
Figure~\ref{fig:example-script} contains several examples.  Each parameter of
the function should be given a type.  

A return type must be given for a recursive function, or if a forward
reference is made to the function: this is because the type checker considers
declarations in the order they appear, but can use claimed types of later
functions.  For example, the return type for |f| is necessary in the
following, since the use of~|f| precedes its definition:
%
\begin{scala}
val y = f(3); def f(x: Int): Int = x+1
\end{scala}
%
However, no return type is necessary in the following, since the use of~|f| is
after its definition:
%
\begin{scala}
def f(x: Int) = x+1; val y = f(3) 
\end{scala}

A function can be defined that takes several vectors of parameters.  For
example: 
%
\begin{scala}
def add(x: Int, y: Int)(z: Int) = x+y+z
val a = add(4,5) 6; val f = add(6,7); val b = f 8
\end{scala}
%
Here |f| is a function of type |Int => Int|.

Parametric polymorphism is supported in a style similar to as in Scala.  Type
parameters are introduced by enclosing them in square brackets after the name
of a function, for example:
%
\begin{scala}
def apply[A,B](f: A => B, x: A): B = f(x)
\end{scala}
%
This defines such a function for all types~|A| and~|B|.  For example
\begin{scala}
def double(x: Float) = 2*x
val a = apply(head, [1,2,3]); val b = apply(double, 3.4)
\end{scala}

Higher-order functions are supported in a style similar to Scala and Haskell. 
For example, we can define a function~|map| such that |map f| applies~|f| to
each member of a list (of the approriate type):
%
\begin{scala}
def map[A,B](f: A => B)(xs: List[A]): List[B] = 
  if(isEmpty xs) [] else f(head xs) :: map f (tail xs)
\end{scala}

Types that allow equality tests, using |==| and |!=|, are called
\emph{equality types}; examples include the atomic types, |List[A]| if |A| is
itself an equality type, but not function types.  Some polymorphic functions
can only be applied to equality types.  This is denoted by including a type
assumption such as \SCALA{A <: Eq} to specify that |A| must be an instance of
the equality type class~|Eq|.  For example, the following function removes the
first occurrence (if any) of a value~|x| from a list; but it uses an equality
test involving~|x|, so the type of~|x| must be an equality type.
%
\begin{scala}
def remove[A <: Eq](x: A)(xs: List[A]): List[A] =
  if(isEmpty xs) [] else if(head xs == x) tail xs else x :: remove x (tail xs)
\end{scala}

The language allows overloading of names of functions declared using |def|s
(but not |val|s).  For example,
\begin{scala}
def double(n: Int) = 2*n
def double(x: Float): 2*x
\end{scala}
%
When an overloaded function is applied, the type checker needs to be able to
determine which instance is intended, for example,
\begin{scala}
val m = double(3) // Int => Int instance.
\end{scala}
%
In some cases, it is necessary to give an explicit type to the argument,
particularly in the case of a cell read; for example,
\begin{scala}
val y = double(#A4: Float) // Float => Float instance.
\end{scala}
%
To allow type deductions to be made, overloaded functions must have disjoint
types for their first vectors of parameters, and may not be polymorphic. 

\subsection{Assertions}

The basic form for an assertion is |assert(|\syntax{expression}|)|, where
the expression should evaluate to a boolean.  When the script is executed, if
that boolean is|false|, an error is generated.

A generalisation is |assert(|\syntax{expression}|, |\syntax{expression}|)|,
where the second expression should evaluate to a string.  If the assertion
fails, that string is included in the error message. 

%%%%%%%%%%%%%%%%%%%%%%%%%%%%%%%%%%%%%%%%%%%%%%%%%%%%%%%

\subsection{Cell writes}
\label{ssec:cell-writes}

The syntax for a cell write is \syntax{cell-expression} = \syntax{expression}.
Figure~\ref{fig:example-script} contains examples where the cell expression on
the left-hand side uses the |Cell| construction.  The \SCALA{\#} construction
can also be used on the left-hand side; for example,
\begin{scala}
#A4 = #A3:Int + 2
\end{scala}

In most simple applications, cell writes occur only at the top level of a
script.  However, they may also appear within a function definition or a
block expression.  For example, the following code is equivalent to the |for|
loop at line~\ref{line:write-facts} of Figure~\ref{fig:example-script}.
%
\begin{scala}
def writeFacts(r: Row): Unit = 
  if(r == firstEmpty) ()
  else{ Cell(out, r) = fact(Cell(In,r)); writeFacts(r+1) }

writeFacts(#0)
\end{scala}
%
The function |writeFacts(r)| writes the factorials into rows \SCALA{r until
  firstEmpty}.  This is invoked via a function call; see
Section~\ref{sec:function-calls}. 

Being able to preform cell writes within a function is useful in applications
where the number of cells written depends upon earlier calculations.  It
allows a recursive function that keeps calculating and writing cells until a
suitable end state is reached.

Recall that the application treats any attempt to write to a non-empty cell as
an error. 

%%%%%%%%%%%%%%%%%%%%%%%%%%%%%%%%%%%%%%%%%%%%%%%%%%%%%%%

\subsection{{\scalashape for} loops}
\label{ssec:for-loops}

The basic form of a |for| loop is
\begin{scala}
for (£\syntax{name}£ <- £\syntax{list}£) £\syntax{statements}£
\end{scala}
where \syntax{statements} is either a single statement (a declaration, cell
write, or |for| loop), or several statements enclosed in braces.  For example,
%
\begin{scala}
for(r <- #0 to #5){ 
  val x = Cell(#A, r): Int; Cell(#B, r) = 2+x; Cell(#C, r) = 2*x
}
\end{scala}

A |for| loop may have multiple generators, or boolean filters; for example,
%
\begin{scala}
for(r <- #0 to #5; c <- [#A, #C]; if r != #3) 
  Cell(c+1,r) = Cell(c,r):Int + 3
\end{scala}

A generator may pattern match against a tuple, for example,
\begin{scala}
for ((r,n) <- [(#1,2), (#2,4), (#3,6)]) 
  Cell(#A, r) = factorial(n)
\end{scala}

Like cell writes, a |for| loop normally occurs just at the top level of a
script, but may also occur within a function definition or a block expression.

%%%%%%%%%%%%%%%%%%%%%%%%%%%%%%%%%%%%%%%%%%%%%%%%%%%%%%%

\subsection{Function calls}
\label{sec:function-calls}

If a function returns the |Unit| type, a function call can be used as a
statement.  Section~\ref{ssec:cell-writes} contained an example. 

%%%%%%%%%%%%%%%%%%%%%%%%%%%%%%%%%%%%%%%%%%%%%%%%%%%%%%%

\subsection{Conditional statements}

Statements can be executed conditionally, using an |IF| statement such as
\begin{scala}
IF(#A3:Int > 3){ #B3 = 4; #B4 = true } ELSE #B5 = false
\end{scala}
The |ELSE| part is optional.  
Note that the keywords |IF| and |ELSE| are capitalised, to distinguish this
statement from a conditional expression.  Each branch may be a single
statement or several statements enclosed in braces.  Note that any declaration
in a branch is local to that branch. 

%%%%%%%%%%%%%%%%%%%%%%%%%%%%%%%%%%%%%%%%%%%%%%%%%%%%%%%

\subsection{{\scalashape \#include} statements}

A script may import the contents of another script via an \SCALA{#include}
statement, for example
\begin{scala}
#include "other-file.dir"
\end{scala}
The name of the other file should be enclosed in quotation marks.  There
should be nothing following on the same line.
